

\subsection{Robustness and Uncertainty}
\label{sec:robustness}


Many researches have shown that deep models are easily fooled by small and unrecognizable perturbations on the input images, a phenomenon referred to as adversarial attacks~\cite{fgsm, szegedy2013intriguing}.
One straightforward way to enhance robustness and uncertainty is an input augmentation by generating unseen samples~\cite{madry2017towards}.
We evaluate robustness and uncertainty improvements due to input augmentation methods including Mixup, Cutout, and CutMix.

\noindent\textbf{Robustness: }
We evaluate the robustness of the trained models to adversarial samples, occluded samples, and in-between class samples.
We use ImageNet pre-trained ResNet-50 models with same setting as in Section~\ref{sec:imagenet_cls}.

Fast Gradient Sign Method (FGSM)~\cite{fgsm} is used to generate adversarial perturbations and we assume that the adversary has full information of the models (\textit{white-box attack}).
We report top-1 accuracies after attack on ImageNet validation set in Table~\ref{table:fgsm}. 
CutMix significantly improves the robustness to adversarial attacks compared to other augmentation methods.

For occlusion experiments, we generate occluded samples in two ways: center occlusion by filling zeros in a center hole and boundary occlusion by filling zeros outside of the hole.
In Figure~\ref{fig:robustness_occ}, we measure the top-1 error by varying the hole size from $0$ to $224$. For both occlusion scenarios, Cutout and CutMix achieve significant improvements in robustness while Mixup only marginally improves it.
Interestingly, CutMix almost achieves a comparable performance as Cutout even though CutMix has not observed any occluded sample during training unlike Cutout.

Finally, we evaluate the top-1 error of Mixup and CutMix in-between samples. 
The probability to predict neither two classes by varying the combination ratio $\lambda$ is illustrated in Figure~\ref{fig:robustness_inbetween}.
We randomly select $50,000$ in-between samples in ImageNet validation set.
In both experiments, Mixup and CutMix improve the performance while improvements due to Cutout are almost negligible.
Similarly to the previous occlusion experiments, CutMix even improves the robustness to the unseen Mixup in-between class samples.


\begin{table}[]
% \small
\centering
\begin{tabular}{@{}lcccc@{}}
\toprule
               & Baseline & Mixup & Cutout & CutMix        \\ \midrule
Top-1 Acc (\%) & 8.2      & 24.4  & 11.5   & \textbf{31.0} \\ \bottomrule
\end{tabular}
\caption{Top-1 accuracy after FGSM white-box attack on ImageNet validation set.}
\label{table:fgsm}
\end{table}

\begin{table}[]
% \small
\tabcolsep=0.1cm
\centering
\begin{tabular}{@{}lccc@{}}
\toprule
Method   & TNR at TPR 95\%              & AUROC                       & Detection Acc. \\ \midrule
Baseline & 26.3 (+0)                    & 87.3 (+0)                   & 82.0 (+0)                    \\
Mixup    & 11.8 \redp{-14.5}            & 49.3 \redp{-38.0}           & 60.9 \redp{-21.0}           \\
Cutout   & 18.8 \redp{-7.5}            & 68.7 \redp{-18.6}            & 71.3 \redp{-10.7}            \\
CutMix   & \textbf{69.0 \greenp{+42.7}} & \textbf{94.4 \greenp{+7.1}} & \textbf{89.1 \greenp{+7.1}} \\ \bottomrule
\end{tabular}
\vspace{-0.15cm}
\caption{Out-of-distribution (OOD) detection results with CIFAR-100 trained models. Results are averaged on seven datasets. All numbers are in percents; higher is better. 
}
\vspace{-0.3cm}
\label{table:OOD}
\end{table}


\noindent\textbf{Uncertainty: }
We measure the performance of the out-of-distribution (OOD) detectors proposed by~\cite{hendrycks2016baseline} which determines whether the sample is in- or out-of-distribution by score thresholding.
We use PyramidNet-200 trained on CIFAR-100 datasets with same setting as in Section~\ref{sec:CIFAR-100-classification}.
In Table~\ref{table:OOD}, we report the averaged OOD detection performances against seven out-of-distribution samples from~\cite{hendrycks2016baseline, liang2017odin}, including TinyImageNet, LSUN~\cite{yu2015lsun}, uniform noise, Gaussian noise, etc. More results are illustrated in Appendix~\ref{appendix:robustness}.
Mixup and Cutout augmentations aggravate the over-confidence of the base networks. Meanwhile, CutMix significantly alleviates the over-confidence of the model.

